\chapter{Sleeve Gastrectomy}
\section*{Overview}
Sleeve gastrectomy is bariatric surgery that removes most of the stomach, leaving a narrow sleeve-shaped stomach.
\section*{Common Instruments and Equipment}
The exact setup varies by technique and setting. Common items may include:
\begin{itemize}[nosep,leftmargin=*]
  \item laparoscopic camera and ports
  \item long bariatric stapler and reloads
  \item gastric sizing bougie and leak-test equipment when used
\end{itemize}

\section*{Why It Is Done}
It is used for selected patients with severe obesity and related disease after comprehensive medical and lifestyle evaluation.
\section*{How It Is Performed}
Usually laparoscopically, the surgeon removes the greater portion of the stomach along a calibration tube and closes the remaining stomach with a staple line.
\section*{Contraindications}
No single contraindication applies to every patient; the treating team assesses whether the procedure is suitable for the individual and the reason for treatment.

\section*{Precautions}
The care team reviews relevant health conditions, medicines, allergies, and procedure-specific risks, and explains any changes needed for the individual.

\section*{Risks and Recovery}
Bleeding, leak, infection, reflux, stricture, blood clots, and nutritional deficiencies are possible. Long-term dietary, vitamin, laboratory, and weight-management follow-up is required.
\section*{References}
\begin{enumerate}\item MedlinePlus. Bariatric Surgery. U.S. National Library of Medicine; 2025.\item Current Medical Diagnosis and Treatment. McGraw-Hill; 2025.\end{enumerate}
\clearpage
